% part: intuitionistic-logic
% chapter: propositions-as-types
% section: introduction

\documentclass[../../../include/open-logic-section]{subfiles}

\begin{document}

\olfileid{int}{pty}{int}

\olsection{پېژندنه}

لامبډا حساب د محاسبې يو ساده ماډل دے۔ دا له متغيرونو جوړو
ترمونو سره کار کوي۔ که $N$ او $M$ ترمونه وي، نو $(NM)$ هم
ترم دے («په~$M$ د $N$ تطبيق»)۔ که $N$ ترم وي، نو
$\lambd[x][N]$ هم ترم دے۔ که $N$ متغير~$x$ ولري، نو په
$\lambd[x][N]$ کښې د~$x$ اړوند آزاد موردونه د $\lambd[x]$
عملګر تړي، او بيا آزاد نۀ پاتې کېږي۔ د $(\lambd[x][N]M)$
په بڼه ترم $\Subst{N}{M}{x}$ ته $\beta$-انقباض کوي
(په~$N$ کښې د $x$ د هر آزاد مورد پر ځاے د $M$ د اېښودلو
پايله؛ د آزاد متغير د نيولو د مخنيوي دپاره تړلي نومونه د اړتيا
پر وخت مخکې بيا نومول کېږي)۔ ترم $N$ په يوه ګام کښې~$N'$
ته $\beta$-بدلون کوي که $N'$ په $N$ کښې د يوه فرعي ترم
په $\beta$-انقباض تر لاسه شي؛ $N \redone N'$ ليکو۔ په
$\lambd$-حساب کښې محاسبه د $N \redone N_1
\redone N_2 \redone \dots$ په بڼه لړۍ ده۔ که دغه لړۍ په
داسې ترم $N_k$ پاے ته ورسي چې هېڅ فرعي ترم ئې
$\beta$-انقباض نۀ شي کولے، نو دغه ترم \emph{عادي} يا
د~$N$ عادي بڼه بلل کېږي۔ په $\lambd$-حساب کښې محاسبه د
همداسې لړيو په توګه وګڼئ۔ محاسبه هغه وخت درېږي چې عادي بڼې
ته ورسي، او همدغه عادي بڼه د محاسبې پايله ده۔ د محاسبې دغه
ساده ماډل د ټيورينګ له پلوه بشپړ دے: هره جزوي بازګشتي تابع په
$\lambd$-حساب کښې د يوه ترم په وسيله محاسبه کېدے شي۔

هاسکل کَري او ويليم هاورډ په خپلواک ډول د فطري استنتاج او
$\lambd$-حساب تر منځ نږدې ورته‌والے وموند۔ په ابتدايي معنا،
ترم $\lambd[x][N]$ يوه تابع ده: $x$ ئې آرګومېنټ دے؛ $N$
ښيي چې د ورکړل شوي~$x$ دپاره ارزښت څنګه محاسبه کړو۔ نو
ترم $(\lambd[x][N]M)$ په آرګومېنټ~$M$ د تابعې
$\lambd[x][N]$ تطبيق دے؛ $\beta$-انقباض ئې د ارزښت د
موندلو طريقه ښيي: په $N$ کښې د $x$ د آزادو موردونو پر ځاے
$M$ کېږدئ (او د پايلې ترم ارزونه جاري وساتئ)۔ اوس دا د
ټاکلي داخلې او وتلو سېټ لرونکې تابعې وګڼئ۔ که د
$\lambd[x][N]$ د داخلې سېټ $A$ او د وتلو سېټ~$B$ وي،
نو متغير~$x$ يوازې له~$A$ څخه آرګومېنټونه اخلي او $N$
له~$B$ څخه ارزښتونه ورکوي۔ له دې کبله، $(\lambd[x][N]M)$
يوازې هغه وخت معنا لري چې $\lambd[x][N]$ د $A \to B$
تابع او $M \in A$ وي۔ په $\lambd$-حساب کښې د دغو
معلوماتو په ورزياتولو \emph{ټايپ‌لرونکے} $\lambd$-حساب
تر لاسه کېږي۔ په ټايپ‌لرونکي $\lambd$-حساب کښې هره څرګندونه
$(\lambd[x][N]M)$ روا نۀ ده؛ يوازې هغه روا ده چې د
$\lambd[x][N]$ او $M$ «ټايپونه» سره سمون خوري، يعنې
$\lambd[x][N]$ ټايپ $A \lif B$ او $M$ ټايپ~$A$ ولري۔
د $\lambd[x][N]$ ټايپ يوازې هغه وخت $A \to B$ وي چې
د $x$ ټايپ $A$ او د $N$ ټايپ~$B$ وي۔ په عمومي ډول، هره
څرګندونه $(M_1M_2)$ روا نۀ ده۔ يوازې هغه ترمونه $(M_1M_2)$
روا دي چې $M_1$ ټايپ $A \to B$ او $M_2$ ټايپ~$A$
ولري؛ په دغه صورت کښې د $(M_1M_2)$ ټايپ~$B$ دے۔

د ټايپ ټاکلو دغه قواعد اوس په څرګند ډول د فطري استنتاج له
!!{derivation}s سره سمون لري: ټايپونه د !!{formula}s په وسيله
ښودل کېږي، او خپله !!{derivation}s له $\lambd$-ترمونو سره
سمون لري۔ د بېلګې په توګه، د $!A \lif !B$ !!a{derivation}
له داسې ترم $\lambd[\typeof{x}{!A}][N]$ سره سمون لري چې
د تابعې ټايپ ئې $!A \lif !B$ دے۔ د فطري استنتاج استنتاجي
قواعد د ټايپ‌لرونکي لامبډا حساب د ټايپ ټاکلو له قواعدو سره
سمون لري؛ د بېلګې په توګه،
\begin{prooftree}
  \AxiomC{$\Discharge{!A}{x}$}
  \DeduceC{$!B$}
  \DischargeRule{\Intro\lif}{x}
  \UnaryInfC{$!A \lif !B$}
  \DisplayProof\bottomAlignProof
  \quad\text{سره سمون لري}\quad
  \Axiom$x: !A \fCenter N: !B$
  \UnaryInf$ \fCenter \lambd[\typeof{x}{!A}][N] : !A \lif !B$
\end{prooftree}
د ښي اړخ قاعده وايي چې که $x$ ټايپ~$!A$ او $N$
ټايپ~$!B$ ولري، نو $\lambd[\typeof{x}{!A}][N]$
ټايپ~$!A \lif !B$ لري۔

د $\Elim\lif$ قاعده د تطبيق د ترمونو د ټايپ ټاکلو له
قاعدې سره سمون لري، يعنې
\[
  \AxiomC{$!A \lif !B$}
  \AxiomC{$!A$}
  \RightLabel{\Elim\lif}
  \BinaryInfC{$!B$}
  \DisplayProof
  \quad\text{سره سمون لري}\quad
  \Axiom$\fCenter M_1: !A \lif !B$
  \Axiom$\fCenter M_2: !A$
  \BinaryInf$ \fCenter (M_1M_2) : !B$
  \DisplayProof
\]
که د \Intro\lif{} قاعدې سمدستي ورپسې د \Elim\lif{}
قاعده راشي، نو !!{derivation} ساده کېدے شي:
\begin{gather*}
  \AxiomC{$\Discharge{!A}{x}$}
  \DeduceC{$!B$}
  \DischargeRule{\Intro\lif}{x}
  \UnaryInfC{$!A \lif !B$}
  \AxiomC{}
  \DeduceC{$!A$}
  \RightLabel{\Elim\lif}
  \BinaryInfC{$!B$}
  \DisplayProof
  \quad\redone\quad
  \AxiomC{}
  \DeduceC{$!A$}
  \DeduceC{$!B$}
  \DisplayProof
\end{gather*}
دا د لامبډا ترمونو له $\beta$-انقباض سره سمون لري:
\[
(\lambd[\typeof{x}{!A}][N]M) \quad\redone\quad
\Subst{N}{M}{x}.
\]

د $\land$ د قواعدو او د «ضرب» د ټايپونو تر منځ، او د
$\lor$ د قواعدو او د «جمع» د ټايپونو تر منځ هم ورته سمونونه شته۔

په ساده ټايپ‌لرونکي لامبډا حساب کښې د ترمونو او د فطري
استنتاج د !!{derivation}s دغه سمون «کَري--هاورډ»، «ثبوتونه د
پروګرامونو په توګه» يا «بيانونه د ټايپونو په توګه» سمون بلل
کېږي۔ د !!{formula}s (بيانونو) له ټايپونو او د ثبوتونو له
ترمونو سره د سمون تر څنګ، سمونونه داسې لنډ ښودلے شو:
\begin{center}
  \begin{tabular}{c c c}
    منطق & پروګرام \\
    \hline
    بيان/!!{formula} & ټايپ \\
    ثبوت & ترم \\
    فرض & متغير \\
    ايستل شوے فرض & تړلے متغير \\
    پاتې پرانيستے فرض & آزاد متغير \\
    $!A \lif !B$ & د تابعې ټايپ \\
    $!A \land !B$ & د ضرب ټايپ \\
    $!A \lor !B$ & د جمع ټايپ \\
    $\lfalse$ & تش ټايپ \\
    \hline
  \end{tabular}
\end{center}

د کَري--هاورډ سمون د اتومات ثبوت مرستيالانو او د پروګرامونو د
ټايپ کتونکو يو بنسټ دے، ځکه د يوه بيان شاهد ثبوت کتل (لکه
پورته مو چې وکړل) د دې کتلو سره برابر دي چې پروګرام (ترم)
خپل اعلان شوے ټايپ لري که نه۔
\end{document}
